\documentclass[11pt,a4paper]{article}
\usepackage[utf8]{inputenc}
\usepackage{amsmath, amsfonts, amssymb, amsthm}
\usepackage{geometry}
\geometry{a4paper, margin=1in}
\usepackage{hyperref}
\usepackage{braket}

\title{Finite-Dimensional Fibonacci-Peirce Code Algebra and its Level-Two Cuntz Transport}
\author{Autonomously Formalized via Lean 4}
\date{\today}

\newtheorem{theorem}{Theorem}[section]
\newtheorem{lemma}[theorem]{Lemma}
\newtheorem{definition}[theorem]{Definition}

\begin{document}

\maketitle

\begin{abstract}
We construct a finite-dimensional Fibonacci-Peirce code algebra, prove the Knill-Laflamme compression condition for a localized family of errors, and mathematically transport this structure into a level-two finite core generated by a Cuntz family. The entire construction is formalized and verified in the Lean 4 theorem prover.
\end{abstract}

\section{Algebraic Setup}
We define a finite-dimensional algebraic core generated by a Cuntz family $S_1, S_2$ satisfying $S_i^* S_j = \delta_{ij}$ and $S_1S_1^* + S_2S_2^* = 1$. For words of length two, $\mu = (i,j)$, we define length-two words and their corresponding transitions:
\begin{equation}
S_\mu = S_i S_j, \qquad e_{\mu\nu} = S_\mu S_\nu^*.
\end{equation}

\section{Fibonacci-Peirce Projector}
We introduce the fundamental Fibonacci matrix:
\begin{equation}
M_F =
\begin{pmatrix}
0 & 1 \\
1 & 1
\end{pmatrix}, \qquad M_F^2 = M_F + I.
\end{equation}
By isolating the dominant eigenvector $\ket{v_+}$, we construct the corresponding Peirce spectral projector:
\begin{equation}
P_+ = \ket{v_+}\bra{v_+}.
\end{equation}

\section{Tensor-Amplified Code}
To introduce non-trivial error correction capabilities, we amplify the system to a bipartite tensor product space. We define the code subspace projector as:
\begin{equation}
P_C = P_+ \otimes I_2.
\end{equation}

\section{Level-Two Cuntz Corner}
We construct a non-unital $*$-homomorphism $\Phi_2 : M_4(\mathbb{C}) \to A$ from the algebra of $4 \times 4$ complex matrices to the Cuntz word-two core:
\begin{equation}
\Phi_2(A) = \sum_{\mu,\nu} A_{\mu\nu} e_{\mu\nu}.
\end{equation}
The unit of this level-two corner is precisely the mapped identity matrix, proving that $\Phi_2(I_4) = \widehat{I_4}$.

\section{Logical Matrix Units}
The logical matrix units for our encoded qubit are constructed as $L_{ab} = P_+ \otimes E_{ab}$. Transporting these to the Cuntz algebra gives:
\begin{equation}
\widehat{L}_{ab} = \Phi_2(L_{ab}).
\end{equation}
We verify the fundamental properties of these units as a supported matrix-unit system. In particular, only the diagonal elements $\widehat{L}_{aa}$ are self-adjoint, while $\widehat{L}_{ab}^* = \widehat{L}_{ba}$ forms the involution relation.

\section{Knill-Laflamme Transport}
We consider an error family $E_\alpha = F_\alpha \otimes I_2$ localized to the first tensor factor. Over the logical matrix units, we verify the scalar compression:
\begin{equation}
L_{rs} E_\alpha^* E_\beta L_{tu} = \lambda_{\alpha\beta} \delta_{st} L_{ru},
\end{equation}
where $\lambda_{\alpha\beta} = \braket{v_+ | F_\alpha^* F_\beta | v_+}$. Through the Cuntz map $\Phi_2$, this exact relation is preserved:
\begin{equation}
\widehat{L}_{rs} \widehat{E}_\alpha^* \widehat{E}_\beta \widehat{L}_{tu} = \lambda_{\alpha\beta} \delta_{st} \widehat{L}_{ru}.
\end{equation}
This forms our main algebraic result: a nontrivial logical $M_2(\mathbb{C})$ corner with transported KL relations.

\section{Logical Wheeler-DeWitt Constraint}
To embed geometric minisuperspace dynamics inside the protected code space, we utilize the algebraic Clifford structure formed by the logical matrix units. We define the logical Pauli operators:
\begin{equation}
Z_L = L_{00} - L_{11}, \qquad X_L = L_{01} + L_{10}.
\end{equation}
These generators satisfy the algebra $Z_L^2 = P_C$, $X_L^2 = P_C$, and anticommute $\{Z_L, X_L\} = 0$. Using these, we construct the logical Wheeler-DeWitt Hamiltonian parameterized by coordinates $(t,x,y) \in \mathbb{R}^3$:
\begin{equation}
H(t,x,y) = t P_C + x Z_L + y X_L,
\end{equation}
and its algebraic conjugate $H^\sharp(t,x,y) = t P_C - x Z_L - y X_L$. The geometry of the (1,2) Minkowski minisuperspace emerges exactly as the null constraint of this operator inside the logical code corner:
\begin{equation}
H H^\sharp = (t^2 - x^2 - y^2) P_C.
\end{equation}
Thus, the algebraic Wheeler-DeWitt zero-energy condition $H H^\sharp = 0$ is natively resolved inside the $P_C$ coherent subspace.

\subsection{Null-Cone Factorization and Nonvanishing Criteria}
The algebraic factorization dictates that on the null-cone $t^2 - x^2 - y^2 = 0$, we have $H H^\sharp = H^\sharp H = 0$. In the abstract $*$-algebra, the nonvanishing of the operators themselves relies on the injectivity of the logical embedding map. We rigorously prove the following theorem-honest classification:
\begin{quote}
Every injective realization of the logical $M_2(\mathbb{C})$-corner transports nontrivial null matrices to non-zero mutually annihilating elements.
\end{quote}
Thus, provided $\operatorname{Function.Injective}(\texttt{logicalMatrixEmbedding})$, the condition $(t,x,y) \neq (0,0,0)$ on the null-cone yields the nontrivial zero-divisor pair $H \neq 0$ and $H^\sharp \neq 0$.

\subsection{Zorn Minisuperspace Connection}
To connect this structure back to the global geometry of the 8D split-quaternion/Zorn layer, the matrix formulation $H_{\mathrm{mat}}(t,x,y)$ of the Wheeler-DeWitt Hamiltonian is embedded into the Zorn algebra. The determinant is evaluated exactly as $\det H_{\mathrm{mat}} = t^2 - x^2 - y^2$. Consequently, the Zorn norm transport explicitly bridges the Clifford corner back to the macroscopic spacetime geometry via:
\begin{equation}
N_{\mathrm{Zorn}} \bigl(\iota_e(H_{\mathrm{mat}}(t,x,y))\bigr) = t^2 - x^2 - y^2.
\end{equation}
This corollary machine-verifiably locks the microscopic logical Clifford factor to the macroscopic Zorn spacetime.

\section{Logical Hodge Laplacian and Spectral Splitting}

Let
\[
\{L_{ab}\}_{a,b\in\{0,1\}}
\]
be the supported logical matrix units constructed in the level-two Cuntz core. They satisfy
\[
L_{ab}L_{cd}=\delta_{bc}L_{ad},
\qquad
L_{ab}^{*}=L_{ba},
\qquad
L_{00}+L_{11}=P_C,
\]
where $P_C$ is the support projection of the logical corner.

\subsection{Logical grading}

Define the logical grading element
\[
\Gamma:=L_{00}-L_{11}.
\]
The matrix-unit relations imply
\[
\Gamma^{*}=\Gamma,
\qquad
\Gamma^{2}=P_C.
\]
Thus $\Gamma$ is a self-adjoint involution relative to the supported corner $P_C A P_C$.

\subsection{Inner derivation}

The grading determines the inner derivation
\[
\delta_{\Gamma}(A)
:=
[\Gamma,A]
=
\Gamma A-A\Gamma.
\]
It satisfies the Leibniz identity
\[
\delta_{\Gamma}(AB)
=
\delta_{\Gamma}(A)B
+
A\delta_{\Gamma}(B).
\]

We define the associated logical Hodge Laplacian by the double commutator
\[
\mathcal H_{\Gamma}(A)
:=
\delta_{\Gamma}^{2}(A)
=
[\Gamma,[\Gamma,A]].
\]

For every element $A$ supported by $P_C$, namely $P_C A = A = A P_C$, one has
\[
\mathcal H_{\Gamma}(A) = 2\bigl(A-\Gamma A\Gamma\bigr).
\]

\begin{theorem}[Logical Hodge spectrum]
The logical matrix units satisfy
\[
\mathcal H_{\Gamma}(L_{00})=0,
\qquad
\mathcal H_{\Gamma}(L_{11})=0,
\]
and
\[
\mathcal H_{\Gamma}(L_{01})=4L_{01},
\qquad
\mathcal H_{\Gamma}(L_{10})=4L_{10}.
\]
\end{theorem}

\begin{proof}
Using the matrix-unit relations,
\[
\Gamma L_{00}=L_{00}=L_{00}\Gamma,
\qquad
\Gamma L_{11}=-L_{11}=L_{11}\Gamma,
\]
so both diagonal matrix units commute with $\Gamma$.

For the off-diagonal generators,
\[
\Gamma L_{01}=L_{01},
\qquad
L_{01}\Gamma=-L_{01},
\]
and therefore
\[
\delta_{\Gamma}(L_{01})=2L_{01}.
\]
Applying $\delta_{\Gamma}$ once more gives
\[
\mathcal H_{\Gamma}(L_{01})=4L_{01}.
\]
Similarly,
\[
\delta_{\Gamma}(L_{10})=-2L_{10},
\]
and hence
\[
\mathcal H_{\Gamma}(L_{10})=4L_{10}.
\]
\end{proof}

\begin{corollary}
For
\[
A = a_{00}L_{00} + a_{01}L_{01} + a_{10}L_{10} + a_{11}L_{11}
\]
in the logical $M_2(\mathbb{C})$-subalgebra,
\[
\mathcal H_{\Gamma}(A) = 4a_{01}L_{01} + 4a_{10}L_{10}.
\]
Consequently,
\[
\ker\bigl(\mathcal H_{\Gamma}\vert_{M_2^{\mathrm{logical}}}\bigr) = \operatorname{span}_{\mathbb{C}}\{L_{00},L_{11}\},
\]
whereas
\[
\operatorname{im}\bigl(\mathcal H_{\Gamma}\vert_{M_2^{\mathrm{logical}}}\bigr) = \operatorname{span}_{\mathbb{C}}\{L_{01},L_{10}\}.
\]
\end{corollary}

The operator $\mathcal H_\Gamma$ algebraically separates the diagonal logical observables from the off-diagonal coherence operators. Its kernel on the generated logical $M_2(\mathbb{C})$ is precisely the diagonal subalgebra, while the coherence matrix units have eigenvalue $4$. This spectral splitting provides an algebraic candidate for a dephasing generator.

\section{Logical Hodge Laplacian and Spectral Splitting}

Let
\[
\{L_{ab}\}_{a,b\in\{0,1\}}
\]
be the supported logical matrix units constructed in the level-two
Cuntz core. They satisfy
\[
L_{ab}L_{cd}=\delta_{bc}L_{ad},
\qquad
L_{ab}^{*}=L_{ba},
\qquad
L_{00}+L_{11}=P_C,
\]
where \(P_C\) is the support projection of the logical corner.

\subsection{Logical grading}

Define the logical grading element
\[
\Gamma:=L_{00}-L_{11}.
\]
The matrix-unit relations imply
\[
\Gamma^{*}=\Gamma,
\qquad
\Gamma^{2}=P_C.
\]
Thus \(\Gamma\) is a self-adjoint involution relative to the supported
corner \(P_CAP_C\).

\subsection{Inner derivation}

The grading determines the inner derivation
\[
\delta_{\Gamma}(A)
:=
[\Gamma,A]
=
\Gamma A-A\Gamma.
\]
It satisfies the Leibniz identity
\[
\delta_{\Gamma}(AB)
=
\delta_{\Gamma}(A)B
+
A\delta_{\Gamma}(B).
\]

We define the associated logical Hodge Laplacian by the double
commutator
\[
\mathcal H_{\Gamma}(A)
:=
\delta_{\Gamma}^{2}(A)
=
[\Gamma,[\Gamma,A]].
\]

For every element \(A\) supported by \(P_C\), namely
\[
P_CA=A=AP_C,
\]
one has
\[
\mathcal H_{\Gamma}(A)
=
2\bigl(A-\Gamma A\Gamma\bigr).
\]

\begin{theorem}[Logical Hodge spectrum]
The logical matrix units satisfy
\[
\mathcal H_{\Gamma}(L_{00})=0,
\qquad
\mathcal H_{\Gamma}(L_{11})=0,
\]
and
\[
\mathcal H_{\Gamma}(L_{01})=4L_{01},
\qquad
\mathcal H_{\Gamma}(L_{10})=4L_{10}.
\]
\end{theorem}

\begin{proof}
Using the matrix-unit relations,
\[
\Gamma L_{00}=L_{00}=L_{00}\Gamma,
\qquad
\Gamma L_{11}=-L_{11}=L_{11}\Gamma,
\]
so both diagonal matrix units commute with \(\Gamma\).

For the off-diagonal generators,
\[
\Gamma L_{01}=L_{01},
\qquad
L_{01}\Gamma=-L_{01},
\]
and therefore
\[
\delta_{\Gamma}(L_{01})=2L_{01}.
\]
Applying \(\delta_{\Gamma}\) once more gives
\[
\mathcal H_{\Gamma}(L_{01})=4L_{01}.
\]
Similarly,
\[
\delta_{\Gamma}(L_{10})=-2L_{10},
\]
and hence
\[
\mathcal H_{\Gamma}(L_{10})=4L_{10}.
\]
\end{proof}

\begin{corollary}
For
\[
A=
a_{00}L_{00}
+a_{01}L_{01}
+a_{10}L_{10}
+a_{11}L_{11}
\]
in the logical \(M_2(\mathbb C)\)-subalgebra,
\[
\mathcal H_{\Gamma}(A)
=
4a_{01}L_{01}
+
4a_{10}L_{10}.
\]
Consequently,
\[
\ker
\bigl(
\mathcal H_{\Gamma}
\vert_{M_2^{\mathrm{logical}}}
\bigr)
=
\operatorname{span}_{\mathbb C}
\{L_{00},L_{11}\},
\]
whereas
\[
\operatorname{im}
\bigl(
\mathcal H_{\Gamma}
\vert_{M_2^{\mathrm{logical}}}
\bigr)
=
\operatorname{span}_{\mathbb C}
\{L_{01},L_{10}\}.
\]
\end{corollary}

We additionally introduce the even-odd decomposition defined by:
\[
\mathsf E_{\mathrm{diag}}(A) = \frac12(A+\Gamma A\Gamma), \qquad
\mathsf E_{\mathrm{off}}(A) = \frac12(A-\Gamma A\Gamma).
\]
It naturally follows that $A = \mathsf E_{\mathrm{diag}}(A) + \mathsf E_{\mathrm{off}}(A)$, and
\[
\mathcal H_\Gamma(A) = 4 \mathsf E_{\mathrm{off}}(A).
\]

The operator $\mathcal H_\Gamma$ algebraically separates the diagonal logical observables from the off-diagonal coherence operators. Its kernel on the generated logical $M_2(\mathbb C)$ is precisely the diagonal subalgebra, while the coherence matrix units have eigenvalue $4$. This spectral splitting provides an algebraic candidate for a dephasing generator.

\section{Limitations}
This formalized layer has strict logical boundaries:
\begin{itemize}
    \item We do not construct an analytic $C^*$-algebra structure (norms and limits), relying entirely on an abstract complex $*$-algebra satisfying Cuntz relations.
    \item A fully executable logical recovery channel (e.g. via completely positive trace-preserving maps) is not instantiated. We prove localized error compression algebraically, stopping short of constructing CPTP dynamical maps.
    \item The model exclusively handles noise localized to the first tensor factor $F_\alpha \otimes I_2$. General uncorrelated noise on both tensor factors is beyond the scope of this formalization.
\end{itemize}

\section{Formal Verification}
The entire algebraic construction detailed in this paper has been computationally formalized in Lean 4. The mapping is provided below:

\begin{table}[h]
\centering
\resizebox{\textwidth}{!}{
\begin{tabular}{|l|l|l|l|}
\hline
\textbf{Mathematical statement} & \textbf{Lean declaration} & \textbf{File} & \textbf{Prerequisites} \\ \hline
$M_F^2 = M_F + I$ & \texttt{fibonacciMatrix\_fusion} & \texttt{PeirceTensorQEC.lean} & Matrix algebra \\ \hline
$P_+^2 = P_+$ & \texttt{fibonacciPeircePlusComplex\_idempotent} & \texttt{PeirceTensorQEC.lean} & Eigensystems \\ \hline
$\operatorname{rank}(P_C) = 2$ & \texttt{peirceCodeProjector\_rank\_two} & \texttt{PeirceTensorQEC.lean} & Tensor products \\ \hline
$L_{ab}L_{cd} = \delta_{bc}L_{ad}$ & \texttt{logicalMatrixUnit\_mul} & \texttt{CuntzPeirceLogicalQubit.lean} & Supported units \\ \hline
Resolved KL compression & \texttt{cuntzLogicalMatrixUnit\_error\_compression} & \texttt{CuntzPeirceLogicalQubit.lean} & Cuntz corner map \\ \hline
Clifford corner relations & \texttt{logicalSigmaZ\_anticommute\_logicalSigmaX} & \texttt{LogicalCliffordHamiltonian.lean} & Logic Pauli matrices \\ \hline
Hamiltonian/conjugate factorization & \texttt{logicalWDWHamiltonian\_mul\_conjugate} & \texttt{LogicalCliffordHamiltonian.lean} & Clifford algebra \\ \hline
Matrix nonvanishing classification & \texttt{logicalWDWMatrix\_eq\_zero\_iff} & \texttt{LogicalCliffordHamiltonian.lean} & Matrix algebra \\ \hline
Conditional injective transport & \texttt{logicalWDWHamiltonian\_ne\_zero} & \texttt{LogicalCliffordHamiltonian.lean} & Injective transport \\ \hline
Nontrivial null zero-divisor pair & \texttt{logicalWDW\_nontrivial\_zero\_divisor\_pair} & \texttt{LogicalCliffordHamiltonian.lean} & Null-cone $H=0$ \\ \hline
Zorn norm transport & \texttt{zornLineEmbed\_norm\_logicalWDWMatrix} & \texttt{SplitOctonionZornKKS.lean} & Zorn formulation \\ \hline
Algebraic dephasing & \texttt{logicalDecoherence\_sub\_diagonalPart} & \texttt{LogicalDecoherenceSemigroup.lean} & Hodge Laplacian \\ \hline
\end{tabular}
}
\end{table}

\end{document}
